\chapter{LLVM Backend and Co-Optimization}
\label{chap:llvm}

\section{The Role of LLVM in NumLang}
\label{sec:llvm:intro}

While Cranelift provides instant compilation for JIT execution, the LLVM backend (\texttt{src/codegen/llvm\_backend.rs}, implemented using the \texttt{inkwell} Rust bindings) provides industrial-grade optimization for ahead-of-time (AOT) release builds.

By supercompiling high-level recursion and deforestation into clean, canonical SSA MIR, NumLang presents LLVM with intermediate code that is uniquely receptive to low-level microarchitectural optimization:
\begin{itemize}
    \item Intermediate heap structures have already been eliminated by metacomputation.
    \item Loops have been flattened into contiguous memory traversals.
    \item Recurrences have been reduced to linear companion matrix powers.
\end{itemize}

\section{Type-Based Alias Analysis (TBAA) Tree Synthesis}
\label{sec:llvm:tbaa}

To enable LLVM to vectorize loops aggressively without fearing memory aliasing, NumLang synthesizes explicit Type-Based Alias Analysis (TBAA) metadata trees during code generation:
\begin{enumerate}
    \item The compiler emits a root TBAA metadata node representing all NumLang memory.
    \item Struct fields and array slice allocations are assigned disjoint child nodes in the TBAA hierarchy.
    \item Pointer load and store instructions attach \texttt{!tbaa} tags, proving statically to LLVM's alias analysis passes that array buffers never overlap with scalar environment records.
\end{enumerate}

\section{Parameter noalias Attributes}
\label{sec:llvm:noalias}

For all unique pointer parameters (such as arrays passed into specialized stencil kernels), the LLVM backend attaches the \texttt{noalias} parameter attribute. This guarantees that LLVM's auto-vectorizer can generate unrolled AVX2 SIMD instructions without emitting expensive runtime pointer alias checks.

\section{AVX2 SIMD Unrolled Matrix Power Kernels}
\label{sec:llvm:simd}

For recurrence companion matrix systems (such as the coupled Fibonacci matrix power in \texttt{fib\_matrix}), the LLVM backend synthesizes specialized SIMD kernels using x86-64 AVX2 vector instructions:
\begin{itemize}
    \item Companion $2 \times 2$ matrices are packed into 256-bit SIMD registers (\texttt{<4 x i64>}).
    \item Matrix multiplications execute via parallel vector broadcasts and SIMD multiply-accumulate sequences.
    \item Binary exponentiation loops achieve sub-microsecond runtimes on modern hardware.
\end{itemize}

\section{The Inkwell LLVM Code Generation Pipeline}
\label{sec:llvm:backend_impl}

Listing~\ref{lst:llvm_backend_impl} reproduces the primary LLVM codegen pass from \texttt{src/codegen/llvm\_backend.rs}, emitting TBAA trees and vectorization hints:

\begin{lstlisting}[language=Rust, caption={LLVM Code Generation and TBAA Construction (\texttt{src/codegen/llvm\_backend.rs}).}, label={lst:llvm_backend_impl}]
//!
//! Provides ahead-of-time (AOT) compilation of `MirProgram` to native object files (`.obj`)
//! via LLVM (using the `inkwell` crate).
//!
//! When the `llvm-backend` Cargo feature is enabled, this module produces optimized machine code
//! using LLVM's auto-vectorization, loop unrolling, and link-time optimizations.
//! When disabled, it provides stub interfaces that report `LlvmError::BackendDisabled`.

use std::path::Path;
use thiserror::Error;
use crate::mir::lower::MirProgram;

/// LLVM optimization level.
#[derive(Debug, Clone, Copy, PartialEq, Eq, Default)]
pub enum OptLevel {
    O0,
    O1,
    #[default]
    O2,
    O3,
}

impl From<u8> for OptLevel {
    fn from(val: u8) -> Self {
        match val {
            0 => OptLevel::O0,
            1 => OptLevel::O1,
            3 => OptLevel::O3,
            _ => OptLevel::O2,
        }
    }
}

impl std::str::FromStr for OptLevel {
    type Err = String;

    fn from_str(s: &str) -> Result<Self, Self::Err> {
        match s {
            "0" | "o0" | "O0" => Ok(OptLevel::O0),
            "1" | "o1" | "O1" => Ok(OptLevel::O1),
            "2" | "o2" | "O2" => Ok(OptLevel::O2),
            "3" | "o3" | "O3" => Ok(OptLevel::O3),
            other => Err(format!(
                "Invalid optimization level '{}'. Expected 0, 1, 2, or 3.",
                other
            )),
        }
    }
}

/// Errors occurring during LLVM backend compilation.
#[derive(Debug, Error)]
pub enum LlvmError {
    #[error("LLVM backend is not enabled in this build. Rebuild with `--features llvm-backend` to enable LLVM codegen.")]
    BackendDisabled,

    #[error("LLVM target initialization error: {0}")]
    TargetInitError(String),

    #[error("LLVM type lowering error: {0}")]
    TypeLoweringError(String),

    #[error("LLVM codegen error: {0}")]
    CodegenError(String),

    #[error("IO error: {0}")]
    IoError(#[from] std::io::Error),
}

#[cfg(not(feature = "llvm-backend"))]
pub struct LlvmCompiler {
    opt_level: OptLevel,
}

#[cfg(not(feature = "llvm-backend"))]
impl LlvmCompiler {
    pub fn new(opt_level: OptLevel) -> Self {
        Self { opt_level }
    }

    pub fn opt_level(&self) -> OptLevel {
        self.opt_level
    }

    /// Emit textual LLVM IR with TBAA trees, noalias attributes, and loop vectorization metadata.
    pub fn emit_llvm_ir(&self, program: &MirProgram) -> Result<String, LlvmError> {
        emit_text_llvm_ir(program, self.opt_level)
    }

    /// Compile a `MirProgram` to an object file (`.obj`).
    pub fn compile_mir_to_obj(
        &mut self,
        _program: &MirProgram,
        _output_path: &Path,
    ) -> Result<(), LlvmError> {
        Err(LlvmError::BackendDisabled)
    }

    /// Compile a `MirProgram` to in-memory object bytes.
    pub fn compile_mir_to_obj_bytes(
        &mut self,
        _program: &MirProgram,
    ) -> Result<Vec<u8>, LlvmError> {
        Err(LlvmError::BackendDisabled)
    }
}

#[cfg(feature = "llvm-backend")]
pub struct LlvmCompiler {
    context: inkwell::context::Context,
    opt_level: OptLevel,
    struct_fields: std::collections::HashMap<String, Vec<(String, crate::typecheck::types::Type)>>,
}

#[cfg(feature = "llvm-backend")]
impl LlvmCompiler {
    pub fn new(opt_level: OptLevel) -> Self {
        let context = inkwell::context::Context::create();
        Self {
            context,
            opt_level,
            struct_fields: std::collections::HashMap::new(),
        }
    }

    pub fn opt_level(&self) -> OptLevel {
        self.opt_level
    }

    /// Emit textual LLVM IR with TBAA trees, noalias attributes, and loop vectorization metadata.
    pub fn emit_llvm_ir(&self, program: &MirProgram) -> Result<String, LlvmError> {
        emit_text_llvm_ir(program, self.opt_level)
    }

    /// Compile a `MirProgram` to an object file (.obj) on disk.
    pub fn compile_mir_to_obj(
        &mut self,
        program: &MirProgram,
        output_path: &Path,
    ) -> Result<(), LlvmError> {
        let bytes = self.compile_mir_to_obj_bytes(program)?;
        std::fs::write(output_path, bytes)?;
        Ok(())
    }

    /// Compile a `MirProgram` to in-memory object bytes.
    pub fn compile_mir_to_obj_bytes(
        &mut self,
        program: &MirProgram,
    ) -> Result<Vec<u8>, LlvmError> {
        use inkwell::{
            module::{Linkage, Module},
            passes::PassManager,
            targets::{
                CodeModel, FileType, InitializationConfig, RelocMode, Target, TargetMachine,
                TargetTriple,
            },
            OptimizationLevel,
        };

        let module = self.context.create_module("numlang_module");
        let builder = self.context.create_builder();

        // 1. Lower struct definitions
        self.struct_fields.clear();
        for s in &program.structs {
            self.struct_fields.insert(s.name.clone(), s.fields.clone());
        }

        let mut struct_types = std::collections::HashMap::new();
        for s in &program.structs {
            let st = self.context.opaque_struct_type(&s.name);
            struct_types.insert(s.name.clone(), st);
        }

        for s in &program.structs {
            let st = struct_types.get(&s.name).unwrap();
            let mut field_tys = Vec::with_capacity(s.fields.len());
            for (_, field_ty) in &s.fields {
                let llvm_ty = self.type_to_llvm_basic(field_ty, &struct_types)?;
                field_tys.push(llvm_ty);
            }
            st.set_body(&field_tys, false);
        }

        // 2. Declare all functions
        let mut func_vals = std::collections::HashMap::new();
        for func in &program.functions {
            let mut param_tys = Vec::with_capacity(func.params.len());
            for (_, p_ty) in &func.params {
                let llvm_ty = self.type_to_llvm_basic(p_ty, &struct_types)?;
                param_tys.push(llvm_ty.into());
            }

            let fn_type = if func.return_ty == crate::typecheck::types::Type::Void {
                self.context.void_type().fn_type(&param_tys, false)
            } else {
                let ret_llvm = self.type_to_llvm_basic(&func.return_ty, &struct_types)?;
                ret_llvm.fn_type(&param_tys, false)
            };

            let fn_val = module.add_function(&func.name, fn_type, Some(Linkage::External));
            func_vals.insert(func.name.clone(), fn_val);
        }

        // 3. Declare intrinsic / external functions
        self.declare_intrinsics(&module, &mut func_vals);

        // 4. Compile function bodies
        for func in &program.functions {
            self.compile_mir_function(func, &module, &builder, &struct_types, &func_vals)?;
        }

        // 5. Emit Windows entry point `mainCRTStartup` if `main` is present on Windows
        #[cfg(target_os = "windows")]
        if let Some(&main_fn) = func_vals.get("main") {
            self.emit_windows_entry_point(&module, &builder, main_fn);
        }

        // 6. Run LLVM optimization pipeline
        self.run_optimization_passes(&module);

        // 7. Emit object file bytes via TargetMachine
        Target::initialize_x86(&InitializationConfig::default());
        #[cfg(target_os = "windows")]
        let triple_str = "x86_64-pc-windows-msvc";
        #[cfg(target_os = "linux")]
        let triple_str = "x86_64-unknown-linux-gnu";
        #[cfg(target_os = "macos")]
        let triple_str = "x86_64-apple-darwin";
        #[cfg(not(any(target_os = "windows", target_os = "linux", target_os = "macos")))]
        let triple_str = "x86_64-unknown-linux-gnu";
        let triple = TargetTriple::create(triple_str);
        let target = Target::from_triple(&triple)
            .map_err(|e| LlvmError::TargetInitError(e.to_string()))?;

        let opt_llvm = match self.opt_level {
            OptLevel::O0 => OptimizationLevel::None,
            OptLevel::O1 => OptimizationLevel::Less,
            OptLevel::O2 => OptimizationLevel::Default,
            OptLevel::O3 => OptimizationLevel::Aggressive,
        };

        let target_machine = target
            .create_target_machine(
                &triple,
                "x86-64",
                "+avx2",
                opt_llvm,
                RelocMode::Default,
                CodeModel::Default,
            )
            .ok_or_else(|| {
                LlvmError::TargetInitError("Failed to initialize LLVM target machine".to_string())
            })?;

        let memory_buffer = target_machine
            .write_to_memory_buffer(&module, FileType::Object)
            .map_err(|e| LlvmError::CodegenError(e.to_string()))?;


\end{lstlisting}
